\documentclass[10pt]{article}

% ---------- Document Identity (update these only) ----------
\newcommand{\DocStage}{}
\newcommand{\DocTitle}{Your Road to Tier One SOF}
\newcommand{\ProgramName}{182-Week Civilian-to-Selection Readiness Pipeline}
\newcommand{\ProgramSubtitle}{}

% ---------- Page ----------
\usepackage[legalpaper,left=0.5in,right=0.5in,top=0.6in,bottom=0.45in]{geometry}

% ---------- Typography ----------
\usepackage[T1]{fontenc}
\usepackage[utf8]{inputenc}
\usepackage{libertinus}
\usepackage{microtype}
\usepackage{parskip}
\usepackage{ragged2e}
\usepackage{textcomp}
\AtBeginDocument{\color{black}}
\AtBeginDocument{\pagecolor{white}}
\AtBeginDocument{\justifying}

% ---------- Landscape pages ----------
\usepackage{pdflscape}

% ---------- Color ----------
\usepackage[table]{xcolor}
\definecolor{StageBlue}{HTML}{1F4E79}
\definecolor{StageBlueDark}{HTML}{163A5A}
\definecolor{LightGray}{HTML}{F2F4F7}
\definecolor{MidGray}{HTML}{D9DEE5}
\definecolor{RuleGray}{HTML}{C7CED8}
\definecolor{HeaderInk}{HTML}{000000}
\definecolor{Ink}{HTML}{000000}

% ---------- Utility ----------
\usepackage{enumitem}
\sloppy
\setlength{\emergencystretch}{2em}

% ---------- Boxes / UI ----------
\usepackage[most]{tcolorbox}

\tcbset{
  charterTitle/.style={
    enhanced,
    colback=StageBlue!4,
    colframe=StageBlue,
    boxrule=1.25pt,
    arc=3pt,
    left=16pt,right=16pt,top=14pt,bottom=14pt,
    borderline north={3.2pt}{0pt}{StageBlue},
    borderline west={4pt}{0pt}{StageBlue},
    borderline south={1.25pt}{0pt}{StageBlue},
    drop shadow={black!25!white},
  },
  charterRule/.style={
    enhanced,
    colback=LightGray,
    colframe=RuleGray,
    boxrule=0.85pt,
    arc=3pt,
    left=12pt,right=12pt,top=10pt,bottom=10pt,
    borderline west={3pt}{0pt}{StageBlue},
  },
  charterBadge/.style={
    enhanced,
    colback=StageBlue,
    coltext=white,
    colframe=StageBlueDark,
    boxrule=0.6pt,
    arc=2pt,
    left=6pt,right=6pt,top=3pt,bottom=3pt,
  }
}
\newtcbox{\Badge}{nobeforeafter,charterBadge}

% ---------- Lists ----------
\setlist[itemize]{leftmargin=1.5em, itemsep=0.25em, topsep=0.25em}
\setlist[enumerate]{leftmargin=1.8em, itemsep=0.25em, topsep=0.25em}

% ---------- TOC ----------
\usepackage{tocloft}
\setcounter{tocdepth}{3}
\setlength{\cftbeforesecskip}{3pt}
\setlength{\cftbeforesubsecskip}{2pt}
\setlength{\cftbeforesubsubsecskip}{0pt}
\renewcommand{\cftsecfont}{\bfseries}
\renewcommand{\cftsecpagefont}{\bfseries}
\renewcommand{\cftsubsecfont}{\normalfont}
\renewcommand{\cftsubsecpagefont}{\normalfont}
\renewcommand{\cftsubsubsecfont}{\normalfont}
\renewcommand{\cftsubsubsecpagefont}{\normalfont}
\setlength{\cftsecindent}{0pt}
\setlength{\cftsubsecindent}{1.25em}
\setlength{\cftsubsubsecindent}{2.8em}
\setlength{\cftsecnumwidth}{2.1em}
\setlength{\cftsubsecnumwidth}{2.8em}
\setlength{\cftsubsubsecnumwidth}{3.6em}

% Remove the built-in TOC heading so only the custom heading is shown
\makeatletter
\renewcommand{\tableofcontents}{\@starttoc{toc}}
\makeatother

% ---------- Header/Footer: page number only ----------
\usepackage{fancyhdr}
\pagestyle{fancy}
\fancyhf{}
\renewcommand{\headrulewidth}{0pt}
\renewcommand{\footrulewidth}{0pt}
\fancyfoot[C]{\thepage}

% ---------- Section formatting ----------
\usepackage{titlesec}

\titleformat{\section}
  {\Large\bfseries\color{Ink}}
  {\thesection}{0.6em}{}
  [\vspace{0.25em}\textcolor{StageBlue}{\rule{\linewidth}{0.8pt}}]

\titleformat{\subsection}
  {\large\bfseries\color{Ink}}
  {\thesubsection}{0.6em}{}

\titleformat{\subsubsection}
  {\normalsize\bfseries\color{Ink}}
  {\thesubsubsection}{0.6em}{}

\titlespacing*{\section}{0pt}{0.95em}{0.35em}
\titlespacing*{\subsection}{0pt}{0.65em}{0.25em}
\titlespacing*{\subsubsection}{0pt}{0.55em}{0.2em}

% ---------- Table engine ----------
\usepackage{tabularray}
\UseTblrLibrary{booktabs}

% ---------- tabularray: GLOBAL table treatment ----------
\SetTblrInner{
  cells = {fg=black, bg=white, valign=t},
}

% ---------- Header helpers ----------
\newcommand{\Hdr}[1]{\shortstack[c]{\strut #1\strut}}

% ---------- Lines ----------
\newcommand{\TitleLine}{\textcolor{StageBlue}{\rule{0.98\linewidth}{0.95pt}}}
\newcommand{\SubTitleLine}{\textcolor{RuleGray}{\rule{0.98\linewidth}{0.65pt}}}

% ---------- Continued on next page ----------
\newcommand{\ContinuedOnNextPageInline}{%
  \par
  \vfill
  \noindent\textcolor{RuleGray}{\rule{\linewidth}{1pt}}\vspace{-2pt}
  \noindent\hfill\textit{Continued on next page}\par\vspace{-10pt}
  \noindent\textcolor{RuleGray}{\rule{\linewidth}{1pt}}
  \clearpage
}

% ---------- PDF hyperlinks ----------
\usepackage[hidelinks]{hyperref}
\hypersetup{
  colorlinks=false,
  pdfborder={0 0 0},
  linktoc=all,
  pdftitle={\DocTitle},
  pdfauthor={\ProgramName}
}

% =========================================================
\begin{document}

\subsection*{Phase 07 Card Header}
\phantomsection
\addcontentsline{toc}{subsection}{Phase 07 Card Header}

\begin{itemize}
\item Phase \textbf{ID}: Phase 07
\item Phase Name: Advanced Hybrid Overload
\item Duration weeks: 12
\item Version: v1.0
\item Stage Alignment: Stage 5
\item Governing Status: Draft — Non-Deployable
\end{itemize}

\subsection*{1. Phase Mission}
\phantomsection

Consolidate multi-domain readiness under controlled overload by producing repeatable, admissible gate evidence across \textbf{RUN}, \textbf{RUCK}, \textbf{SWIM}, \textbf{CAL}, and \textbf{STR} within protected windows, without violating anti-stacking/spacing intent or compromising durability.

\subsection*{2. Phase Purpose}
\phantomsection

\begin{itemize}
\item Establish Phase 07 as a controlled hybrid verification phase where multi-domain exposure is present but evidence integrity remains the priority constraint.
\item Prepare the pipeline for Phase 08 durability and load tolerance emphasis by proving stable execution under staged multi-domain test windows.
\end{itemize}

\subsection*{3. Entry Criteria}
\phantomsection

\begin{itemize}
\item Entry requires prior phase exit posture to be satisfied with admissible evidence per Stage 1.F and Stage 2 validity.
\item Default posture if evidence is incomplete or invalid: \textbf{HOLD} and reslot to the next protected window per Stage 3, without compressing or stacking.
\item Durability posture at entry must be stable enough that gait-altering pain is not active and foot breakdown is not in escalation posture (\textbf{DURABILITY} constraints apply as gate disqualifier signals).
\end{itemize}

\subsection*{4. Operating Constraints}
\phantomsection

\begin{itemize}
\item Six sessions per week is mandatory.
\item Required session elements per Stage 1.F in correct order:
  \begin{itemize}
  \item Safety self-check first
  \item Warm-up
  \item Mobility
  \item Main work
  \item Cardio
  \item Cooldown
  \item Recovery notes and any required post-session notes per Stage 1.F
  \end{itemize}
\item Cardio minimum standard: minimum 30 minutes continuous per session.
\item 10,000 steps per day is a global requirement and must be logged daily.
\item Validity doctrine: admissible \textbf{VALID} evidence only for gates; invalid tests provide no gate credit and default posture is \textbf{HOLD} with reslot.
\item Safety overrides schedule and performance; stop posture and escalation posture are controlling.
\item Phase 07 does not authorize deviation from Stage 3 protected-window doctrine; gate-grade testing is protected-window only.
\end{itemize}

\subsection*{5. Primary Adaptations and Physiological Focus}
\phantomsection

\begin{itemize}
\item Multi-domain tolerance under controlled fatigue while maintaining evidence validity and safety posture.
\item Improved robustness of \textbf{CAL} battery performance under repeatable standards and artifact posture.
\item \textbf{RUN} verification stability under controlled tool/environment posture and reduced variance near test windows.
\item \textbf{RUCK} capability continuity with strict foot-interface governance and load verification posture.
\item \textbf{SWIM} verification continuity under unit-locked pool verification posture and rest governance.
\item \textbf{STR} proxy stability (submax) as a supporting capability, staged to avoid contaminating \textbf{RUN}/\textbf{RUCK} evidence.
\end{itemize}

\subsection*{6. Primary Modalities}
\phantomsection

\begin{itemize}
\item Emphasized domains: \textbf{RUN}, \textbf{RUCK}, \textbf{CAL}, \textbf{SWIM}, \textbf{STR}, \textbf{DURABILITY}, \textbf{RECOVERY}, \textbf{BODYCOMP} (supporting).
\item De-emphasized domains: \textbf{AER} is optional and supporting only; it must not be used as interchangeable evidence for \textbf{RUN} readiness.
\item Prohibited posture (phase-level): no new tests or constructs outside Stage 2 instruments; no gate-grade tests outside protected windows.
\end{itemize}

\clearpage
\begin{landscape}

\subsection*{7. Weekly Structure}
\phantomsection

\begingroup
\scriptsize
\setlength{\tabcolsep}{2pt}
\renewcommand{\arraystretch}{1.12}

\begin{longtblr}{
  width=\linewidth,
  rowhead=1,
  colspec = {
    Q[l,wd=0.70in]
    X[l,2.05]
    X[l,1.25]
    X[l,2.95]
    X[l,2.90]
  },
  hlines,
  vlines,
  cells = {hsep=6pt, vsep=6pt, halign=l, valign=t},
  cell{1-Z}{1-5} = {fg=black},
  row{odd} = {bg=white},
  row{even} = {bg=LightGray},
  row{1} = {bg=StageBlue!15, fg=black, font=\bfseries, halign=l, valign=t},
}
\Hdr{Session \#} &
\Hdr{Primary Focus} &
\Hdr{Main Work Domains} &
\Hdr{Cardio Modality/Intent} &
\Hdr{Notes/Risk Controls} \\

1 &
Aerobic throughput under stable mechanics &
\textbf{RUN}; \textbf{DURABILITY} &
Modality: step-based locomotion or treadmill as needed; Minutes: 30–60; RPE plus Talk Test: RPE 3–5, full-to-short sentences; Continuity: continuous planned &
Preserve gait quality; enforce footwear stability; downshift immediately on mechanics change. \\

2 &
Calisthenics capacity intent under strict standards &
\textbf{CAL}; \textbf{DURABILITY} &
Modality: low-impact steady; Minutes: 30–45; RPE plus Talk Test: RPE 2–4, full sentences; Continuity: continuous planned &
Avoid shoulder/elbow flare-up drift; \textbf{CAL} battery execution is protected-window only. \\

3 &
Load tolerance continuity and foot integrity governance &
\textbf{RUCK}; \textbf{DURABILITY} &
Modality: low-impact steady; Minutes: 30–60; RPE plus Talk Test: RPE 3–5, full-to-short sentences; Continuity: continuous planned &
Treat foot hotspots as early-warning disqualifier signals; no “push through.” \\

4 &
Strength proxy support and trunk stability &
\textbf{STR}; \textbf{CAL}; \textbf{DURABILITY} &
Modality: low-impact steady; Minutes: 30–45; RPE plus Talk Test: RPE 2–4, full sentences; Continuity: continuous planned &
\textbf{STR} is submax proxy only; no grinding; staged away from maximal \textbf{RUN}/\textbf{RUCK} windows by intent. \\

5 &
Swim technique and verification readiness &
\textbf{SWIM}; \textbf{DURABILITY} &
Modality: low-impact steady; Minutes: 30–45; RPE plus Talk Test: RPE 2–4, full sentences; Continuity: continuous planned &
Pool access and verification posture must be stable; reslot if pool conditions compromise validity. \\

6 &
Consolidation and readiness stabilization &
\textbf{RECOVERY}; \textbf{DURABILITY}; \textbf{BODYCOMP} &
Modality: lowest-risk available; Minutes: 30–45; RPE plus Talk Test: RPE 2–3, full sentences; Continuity: continuous planned &
No novelty near protected windows; preserve recovery markers and documentation completeness. \\

\end{longtblr}

\endgroup

\subsection*{8. Progression Rules}
\phantomsection

\begin{itemize}
\item One progression lever at a time: do not increase multiple primary levers (Cardio duration, complexity, or multi-domain density) in the same week.
\item Preserve intent, reduce risk: regress ROM, leverage, speed, and complexity before reducing frequency (frequency remains 6 sessions/week unless safety stop dictates otherwise).
\item Substitution posture:
  \begin{itemize}
  \item substitute to lower-risk modality to preserve continuous Cardio when foot/surface risk increases,
  \item substitutions that increase risk are prohibited,
  \item all substitutions are logged and tagged per Stage 1.F.
  \end{itemize}
\item Regression triggers:
  \begin{itemize}
  \item gait-altering pain, mechanics change, escalating foot hotspots, sleep collapse driving RPE spikes, repeated \textbf{MODIFIED}/\textbf{INVALID} sessions → downshift and reslot gate attempts.
  \end{itemize}
\item Gate evidence is protected-window only; do not attempt to “prove readiness” in build weeks.
\end{itemize}

\subsection*{9. Deload and Test Placement}
\phantomsection

\begin{itemize}
\item Phase 07 is governed by Stage 3 integrated calendars.
\item Phase 07 protected windows (Stage 3 posture for Phase 06–11 taper-aware pattern):
  \begin{itemize}
  \item Week 1 (baseline) Deload+Test
  \item Week 5 Deload+Test
  \item Week 10 Deload+Test
  \item Week 11 Transition
  \item Week 12 Deload+Test (end gate)
  \end{itemize}
\item Staging posture:
  \begin{itemize}
  \item when multiple domains are assessed in a protected week, the week is treated as a staged battery under Stage 3.A caps/spacing intent,
  \item if spacing cannot be protected, prioritize the decisive gate item and reslot secondary items per Stage 3.E.
  \end{itemize}
\end{itemize}

\end{landscape}

\clearpage
\begin{landscape}

\subsection*{10. Test Battery}
\phantomsection

\begingroup
\scriptsize
\setlength{\tabcolsep}{2pt}
\renewcommand{\arraystretch}{1.12}

\begin{longtblr}{
  width=\linewidth,
  rowhead=1,
  colspec = {
    X[l,1.55]
    X[l,3.05]
    Q[l,wd=0.95in]
    X[l,3.35]
    Q[l,wd=1.35in]
  },
  hlines,
  vlines,
  cells = {hsep=6pt, vsep=6pt, halign=l, valign=t},
  cell{1-Z}{1-5} = {fg=black},
  row{odd} = {bg=white},
  row{even} = {bg=LightGray},
  row{1} = {bg=StageBlue!15, fg=black, font=\bfseries, halign=l, valign=t},
}
\Hdr{Test Timing} &
\Hdr{Test \textbf{ID}/Name} &
\Hdr{Domain} &
\Hdr{Validity Conditions} &
\Hdr{Pass Threshold Stage 2 Tier} \\

Week 5 mid-phase &
\textbf{C17}, only if finalized and deployable, / \textbf{MET-2B-027}-031 Strength proxies &
\textbf{STR} &
Submax posture per protocol; implement and load method logged; warm-up logged; no stop-rule violations; same-day logging complete &
\textbf{INTERMEDIATE} \\

Week 5 mid-phase &
\textbf{C4} / \textbf{MET-2B-010} Push-ups; \textbf{C5} / \textbf{MET-2B-011} Sit-ups; \textbf{C6} / \textbf{MET-2B-012} Pull-ups; \textbf{C7} / \textbf{MET-2B-013} Plank &
\textbf{CAL} &
\textbf{CAL} artifact evidence required for \textbf{VALID} gate use; strict standards per protocol; same-day logging; no stop-rule violations &
\textbf{INTERMEDIATE} \\

Week 10 mid-phase &
\textbf{C10} / \textbf{MET-2B-016} 3-mile run time trial &
\textbf{RUN} &
Measured route or treadmill admissible when settings logged; warm-up logged; environment/tool identity logged; no stop-rule violations; no pauses &
\textbf{INTERMEDIATE} \\

Week 10 mid-phase &
\textbf{C13} / \textbf{MET-2B-022} Swim 500 yd or \textbf{C13} / \textbf{MET-2B-023} Swim 500 m (unit locked per pool) &
\textbf{SWIM} &
No fins; allowable strokes; pool \textbf{ID} and length verification; rest governance honored; same-day logging; no stop-rule violations &
\textbf{INTERMEDIATE} \\

Week 12 end-phase gate &
\textbf{C12} / \textbf{MET-2B-019} 6-mile ruck time trial (35 lb dry load) &
\textbf{RUCK} &
Dry load verified; water logged separately; route/treadmill identity logged; footwear logged; warm-up logged; no stop-rule violations; no pauses &
\textbf{INTERMEDIATE} \\

Week 12 end-phase gate &
\textbf{C4} / \textbf{MET-2B-010}; \textbf{C5} / \textbf{MET-2B-011}; \textbf{C6} / \textbf{MET-2B-012}; \textbf{C7} / \textbf{MET-2B-013} \textbf{CAL} battery &
\textbf{CAL} &
\textbf{CAL} artifact evidence required; strict standards; same-day logging; no stop-rule violations &
\textbf{INTERMEDIATE} \\

\end{longtblr}

\endgroup

\end{landscape}
\clearpage

\subsection*{11. Exit Standards}
\phantomsection

\begin{itemize}
\item Objective gate to advance is based only on admissible evidence produced in protected windows.
\item Intended tier posture alignment for Phase 07 end gate: \textbf{INTERMEDIATE}.
\item Gate requirements:
  \begin{itemize}
  \item End-phase gate window (Week 12) produces \textbf{VALID} evidence for the \textbf{RUCK} 6-mile time trial (\textbf{C12} / \textbf{MET-2B-019}) and a \textbf{VALID} \textbf{CAL} battery (\textbf{C4}–\textbf{C7} / \textbf{MET-2B-010}-013) under protocol conditions.
  \item Any \textbf{INVALID} result provides no gate credit and triggers \textbf{HOLD} with reslot per Stage 3.E; do not compress by stacking into build weeks.
  \end{itemize}
\item Supporting posture:
  \begin{itemize}
  \item \textbf{RUN} and \textbf{SWIM} check-ins (Week 10) are used to confirm readiness and detect drift; they do not override durability constraints or allow unsafe escalation.
  \end{itemize}
\end{itemize}

\subsection*{12. Risk Controls and Stop Rules}
\phantomsection

\begin{itemize}
\item Stage 0.5 and Stage 1.F stop posture and escalation posture are controlling; this Phase Card does not rewrite them.
\item Phase 07 specific risk controls:
  \begin{itemize}
  \item Multi-domain staging: do not cluster maximal \textbf{RUN} and maximal \textbf{RUCK} within short windows; apply Stage 3.A intent and reslot overflow.
  \item Foot integrity governance: hotspots and gait drift are treated as disqualifier signals for load-bearing gate attempts; downshift and reslot rather than push through.
  \item Environment hazards: unsafe footing/heat/air quality/pool constraints default to postpone/substitute posture consistent with Stage 3.E; compromised attempts are invalid for gate use.
  \item \textbf{STR} proxy risk control: submax only; no grinding; no true max attempts; reslot if it would contaminate locomotion gate attempts.
  \end{itemize}
\end{itemize}

\clearpage
\begin{landscape}

\subsection*{13. Required Capabilities and Dependencies}
\phantomsection

\begingroup
\scriptsize
\setlength{\tabcolsep}{2pt}
\renewcommand{\arraystretch}{1.12}

\begin{longtblr}{
  width=\linewidth,
  rowhead=1,
  colspec = {
    X[l,1.85]
    X[l,3.15]
    Q[l,wd=1.25in]
    Q[l,wd=1.25in]
    X[l,3.30]
  },
  hlines,
  vlines,
  cells = {hsep=6pt, vsep=6pt, halign=l, valign=t},
  cell{1-Z}{1-5} = {fg=black},
  row{odd} = {bg=white},
  row{even} = {bg=LightGray},
  row{1} = {bg=StageBlue!15, fg=black, font=\bfseries, halign=l, valign=t},
}
\Hdr{Dependency \textbf{ID}} &
\Hdr{Required Capability Category and Tier} &
\Hdr{Due-By week-in-phase} &
\Hdr{Owner} &
\Hdr{Fail Action} \\

\textbf{DEP-1C-012} &
7.01 Strength and Resistance Training Equipment — Tier appropriate for submax proxy testing &
Week 5 &
Equipment Lead &
If not available/stable: \textbf{HOLD} \textbf{STR} proxy gate evidence; reslot \textbf{C17} to next protected window; do not substitute with non-Stage 2 strength tests. \\

\textbf{DEP-1C-013} &
7.07 Load-Bearing Rucking and Carry Systems — stable dry-load configuration and verification method &
Week 12 &
Equipment Lead &
If not verifiable: \textbf{RUCK} test is not admissible; \textbf{HOLD} and reslot; do not “approximate load.” \\

\textbf{DEP-1C-007} \& \textbf{DEP-1C-008} &
7.11 Footwear Systems and Foot Care — stable footwear interface and hotspot controls &
Week 1 and maintained through Week 12 &
Trainee &
If hotspots/blisters escalate: downshift exposure; ruck gate attempt reslotted; document as durability constraint. \\

\textbf{DEP-1C-017} &
7.17 Aquatic Training and Water Safety Equipment — pool verification posture and safety compliance &
Week 10 &
Coach Lead &
If pool conditions/verification unavailable: \textbf{SWIM} test reslotted; do not replace with proxy evidence for \textbf{SWIM} gate equivalency. \\

\textbf{DEP-1C-006} &
7.18 Testing Timing Measurement and Course-Setup Tools — timing, route/pool verification, logging stability &
Week 1 and maintained &
Coach Lead &
If tool drift or missing fields: test becomes \textbf{INVALID}; \textbf{HOLD} and reslot; do not use invalid evidence for gate decisions. \\

\textbf{DEP-1C-017} &
7.12 Recovery Mobility and Tissue-Care Implements — recovery support to preserve durability &
Week 1 and maintained &
Trainee &
If recovery collapse occurs: convert window to check-in-only; reslot maximal items; preserve safety posture. \\

\end{longtblr}

\endgroup

\subsection*{14. Validity Requirements}
\phantomsection

\begin{itemize}
\item Evidence admissibility is governed by Stage 1.F and Stage 2 validity. Only \textbf{VALID} tests and admissible session evidence can support gate outcomes.
\item \textbf{CAL} gate use requires compliant artifact evidence per Stage 2 protocol posture. Missing or non-compliant artifact makes the \textbf{CAL} test \textbf{INVALID} for gate use.
\item Missing required log fields defaults the test to \textbf{INVALID} for gate use. Invalid tests provide no gate credit.
\item Environment/tool identity and stability posture per Stage 2.E is mandatory for comparability; tool drift near gate windows is treated as validity drift.
\item Invalid or missed gate tests are reslotted to the next protected window per Stage 3.E; do not compress by stacking.
\end{itemize}

\end{landscape}

\subsection*{15. Change Control Notes}
\phantomsection

\begin{itemize}
\item Adjustments allowed without patch (documentation required):
  \begin{itemize}
  \item reslot a missed/invalid test to the next protected window within the phase per Stage 3.E,
  \item swap sequencing within a protected week to preserve spacing intent, without changing the test set.
  \end{itemize}
\item Requires patch:
  \begin{itemize}
  \item changing which tests constitute the mid-phase or end-phase gate battery,
  \item adding new protected windows or extending phase duration,
  \item changing cadence pattern.
  \end{itemize}
\item No change may violate Stage 0 constraints, Stage 1 validity doctrine, Stage 2 instrument set, or Stage 3 protected-window posture.
\end{itemize}

\subsection*{16. Compliance Checklist}
\phantomsection

\begin{itemize}
\item Phase \textbf{ID}, name, and duration match locked phase structure.
\item Operating constraints include: 6 sessions/week; required element order; Cardio minimum 30 minutes continuous; 10,000 steps/day; validity posture.
\item Weekly structure table includes Cardio cell structure (Modality; Minutes; RPE plus Talk Test; Continuity continuous planned) for Sessions 1–6.
\item Deload and test windows are identified and align to Stage 3 taper-aware pattern for Phase 07.
\item Test Battery lists Stage 2 tests and metrics only (C\# and \textbf{MET-2B}-\#\#\#) and limits to major gate-critical tests.
\item Exit Standards align to Test Battery and Stage 2 tier posture (\textbf{INTERMEDIATE}) and require admissible \textbf{VALID} evidence.
\item Dependencies use Stage 7 category IDs and owners; unmet dependencies trigger \textbf{HOLD} and reslot posture.
\item Governing Status remains Draft — Non-Deployable because \textbf{C17} remains unresolved in a gate-relevant test row.
\end{itemize}

\end{document}